\section{问题二：日前预测购电与日内滚动回测模型}

\subsection{问题分析}

问题二需在每天 $0{:}00$ 仅利用此前历史数据及已知日历信息，制定并锁定当天购电计划。与问题一相比，负载与光伏的预测偏差可能造成供电缺口，须按交易时刻电价的 $5$ 倍紧急补购，因此关键在于权衡日前购电成本与缺电风险。

针对负载的星期周期与光伏的短期连续性，分别采用扩展窗口 XGBoost 和近 $7$ 日同刻均值预测，并以历史净负载残差分位修正预测结果。随后沿用问题一的线性规划框架，加入日末留电要求，联合优化购电与储能运行。最后以真实负载、光伏逐时段执行锁定计划，通过储能预留和紧急补购处理缺口，并将实际日末储电量传递至次日，检验连续运行的费用与可行性。

\subsection{模型建立}

\subsubsection{决策变量}

本问沿用问题一的决策变量及约束条件。由于问题二采用连续逐日运行方式，日前规划不再要求日末储电量恢复至 $6000\,\mathrm{kWh}$，而设置 $E_{144}=S^{*}$，其中 $S^{*}$ 为计划日末储电量，其取值在后文参数标定中确定。

\subsubsection{日前点预测}

每天 $0{:}00$ 制定购电计划时，当天真实负载与光伏功率尚未知，因此仅利用决策日前的历史数据预测当天 $144$ 个时段的负载与光伏功率，分别记为 $\hat P^{\mathrm{load}}_{t}$ 和 $\hat P^{\mathrm{pv}}_{t}$。对第 $D$ 日进行预测时，仅使用日期早于 $D$ 的历史数据，并可使用当天已知的时段、星期、月份等日历信息。由于负载与光伏具有不同的时间规律，本文分别建立预测模型。

\textbf{（1）负载预测：扩展窗口 XGBoost。}

对附件 2 一月负载按星期进行方差分解，星期组间平方和占总平方和的 $0.89$；隔周相关系数为 $0.99$，而隔日相关系数仅为 $0.37$，表明负载具有显著的星期周期。相应地，近 $7$ 日同刻均值、昨日同刻和同星期同刻预测的 MAE 分别为 $781$、$650$ 和 $319\,\mathrm{kW}$。其中，同星期同刻方法虽能利用周周期，但难以进一步反映近期负载变化，因此需要同时引入星期信息与近期历史特征。

基于上述特点，本文采用扩展窗口 XGBoost\cite{chen2016xgboost}。模型以每天 $144$ 个时段的负载功率为预测对象，输入时段、星期、月份等日历特征，以及前 $1/2/7$ 日同刻负载、近 $3/7/30$ 日同刻均值和近 $7$ 日同刻标准差。每天 $0{:}00$ 仅利用此前历史样本重新训练，并预测当天负载。

扩展窗口 XGBoost 的负载 MAE 为 $136.42\,\mathrm{kW}$。进一步比较训练窗口，扩展窗口、滚动 $30$ 日和滚动 $60$ 日的 MAE 分别约为 $136$、$163$ 和 $187\,\mathrm{kW}$，对应全年购电费用约为 $1533$、$1624$ 和 $1696$ 万元。说明保留较长历史能够提供更多有效的周周期样本，因此最终采用扩展窗口 XGBoost。

\textbf{（2）光伏预测：近 $7$ 日同刻均值。}

光伏的星期组间平方和占比仅为 $0.0007$，而相邻日期相关系数达到 $0.96$，说明其星期规律较弱，但短期连续性较强。因此，本文采用近 $W$ 日同刻均值预测第 $D$ 日光伏功率：
\begin{equation}
\hat{P}_{D,t}^{\mathrm{pv}}
=
\frac{1}{W}
\sum_{k=1}^{W}
P_{D-k,t}^{\mathrm{pv}},
\qquad
t=1,\ldots,144
\label{eq:q2_pv_forecast}
\end{equation}

利用 $1$ 月 $15$--$31$ 日数据比较不同窗长，$W=5\sim7$ 时 MAE 最低，约为 $92\sim96\,\mathrm{kW}$；$W=1$ 和 $W=21$ 时均增至约 $119\,\mathrm{kW}$。综合预测精度与近期数据代表性，最终取 $W=7$。

进一步比较发现，跨季节同星期同刻均值的 MAE 为 $493\,\mathrm{kW}$，包含星期特征的树模型为 $159.33\,\mathrm{kW}$，均高于近 $7$ 日同刻均值的 $153.40\,\mathrm{kW}$，进一步说明光伏预测更适合利用近期同刻信息。此外，负载与光伏逐时段 Pearson 相关系数的全年均值仅为 $0.031$，且二者时间规律明显不同，因此分别采用扩展窗口 XGBoost 和近 $7$ 日同刻均值进行预测。

把预测方法与调度环节放在同一张图里对照（图~\ref{fig:q2_model_compare}）。

\begin{figure}[!htbp]
  \centering
  \includegraphics[width=0.86\textwidth]{Q2_fig3_model_compare.png}
  \caption{不同预测与控制策略的全年总购电费对照}
  \label{fig:q2_model_compare}
\end{figure}

由图~\ref{fig:q2_model_compare} 可见：负载与光伏都用 $7$ 日均值时总购电费最高，说明 $7$ 日均值法并不高效；XGB 的滚动窗越长费用越低，说明 XGB 能有效利用历史；本文“负载扩展窗 XGB $+$ 光伏 $7$ 日均值”整体费用最低，叠加后文净分位、日末 $2400$ 与 $\rho$ 后为 $1369.54$ 万元，低于固定分位 $0.8$ 的 $1376.52$ 万元。

\subsubsection{净负载分位与日前线性规划}

点预测只能给出“最可能的净负载”，不能表达误差方向。日前计划的关键不是把净负载估准，而是决定为不确定的净负载预留多少计划电。先看单个时段、把储能耦合暂时拿掉：多买 $1\,\mathrm{kWh}$ 计划电的成本是 $\pi$；若缺口 $1\,\mathrm{kWh}$，须按 $5\pi$ 紧急补购，而这 $1\,\mathrm{kWh}$ 若已买进本只需 $\pi$，故少备的边际损失 $C_u=5\pi-\pi=4\pi$；多备 $1\,\mathrm{kWh}$ 已付 $\pi$，用不上则弃光（残值记 $0$），故多备的边际损失 $C_o=\pi$。临界覆盖率为
\begin{equation}
\alpha^{*}=\frac{C_u}{C_u+C_o}=\frac{4\pi}{4\pi+\pi}=0.8,
\label{eq:q2_critical}
\end{equation}
$\pi$ 约去后与电价表无关。储能能让多购电部分留存，真实 $C_o$ 略小于 $\pi$，故 $0.8$ 是单时段近似下的覆盖率。

因此不在负载、光伏上分别取分位，而在净负载上做保守修正，避免把“负荷上偏、光伏下偏”拆成两次平移。记实测净负载、点预测净负载与残差
\begin{equation}
\begin{aligned}
N_t&=P^{\mathrm{load}}_{t}-P^{\mathrm{pv}}_{t},\\
\hat N_t&=\hat P^{\mathrm{load}}_{t}-\hat P^{\mathrm{pv}}_{t},\\
r^N_{d',t}&=N_{d',t}-\hat N_{d',t},\qquad d'<D.
\end{aligned}
\label{eq:q2_netdef}
\end{equation}
假设给定时段 $t$ 的残差在日与日之间可交换，经验 $0.8$ 分位取决策日前全部历史（expanding）：最近 $7$ 天只有 $7$ 个点，$0.8$ 分位接近第 $6/7$ 大顺序统计量、方差过大；按季切开又使每季开头样本更少。于是取
\begin{equation}
\tilde N_t=\hat N_t+Q_{0.8}(r^N_{\cdot,t}),
\label{eq:q2_net_quantile}
\end{equation}
送入日前 LP 时保持光伏点预测不变（夜间归零），把保守净负载加回负载侧
\begin{equation}
\tilde P^{\mathrm{load}}_{t}=\max\bigl(\hat P^{\mathrm{pv}}_{t}+\tilde N_t,0\bigr).
\label{eq:q2_load_back}
\end{equation}
以修正后的预测覆盖全天净负载，在 $0{:}00$ 求解日前线性规划：
\begin{equation}
\begin{aligned}
\min\quad &\sum_{t=1}^{144}\pi_t G_t \\
\mathrm{s.t.}\quad
&G_t+\tilde P^{\mathrm{pv}}_{t}\Delta t+d_t=\tilde P^{\mathrm{load}}_{t}\Delta t+c_t+w_t,\\
&E_t=E_{t-1}+\eta c_t-d_t/\eta,\\
&E_{\min}\le E_t\le E_{\max},\\
&0\le c_t,d_t\le P_{\max}\Delta t,\qquad G_t,w_t\ge0,\\
&E_{144}=S^{*}=2400\,\mathrm{kWh}.
\end{aligned}
\label{eq:q2_lp}
\end{equation}
其中 $E_0$ 取上一日实际执行末值（$2$ 月 $1$ 日为 $6000\,\mathrm{kWh}$）；日末 $2400\,\mathrm{kWh}$ 只约束计划轨迹，实际执行并不重置到该值，日前 LP 只锁死计划购电 $G_t$，充放电按日内回放重算。

\subsubsection{日内 $\rho$ 保护回放}

“保护”指缺口到来时不把储能一次放到下限，“回放”指计划购电锁死后用当天真值逐时段执行。记净残差
\begin{equation}
r_t=P^{\mathrm{load}}_{t}\Delta t-P^{\mathrm{pv}}_{t}\Delta t-G_t.
\label{eq:q2_resid}
\end{equation}
若 $r_t\le0$，仍按贪心充电、剩余弃光，紧急电不用于充电；若 $r_t>0$，并不把储能放到 $E_{\min}$，而是先按日前计划储电量构造预留线
\begin{equation}
\underline{E}_t=E_{\min}+\rho\bigl(E^{\mathrm{plan}}_{t}-E_{\min}\bigr),
\label{eq:q2_reserve}
\end{equation}
再给出放电上限
\begin{equation}
\min\bigl(r_t,\;\eta(E_{t-1}-\underline{E}_t),\;P_{\max}\Delta t\bigr),
\label{eq:q2_dmax}
\end{equation}
其余记紧急购电 $e_t$，按 $5\pi_t$ 结算。$\rho\in[0,1]$ 是预留比例：$\rho=0$ 即原贪心（放到下限 $1200\,\mathrm{kWh}$），$\rho=1$ 则完全按日前计划储电量放电，越大越不提前把次日要用的库存放完。全天总费用
\begin{equation}
C=\sum_{t}\pi_t G_t+\sum_{t}5\pi_t e_t,
\label{eq:q2_cost}
\end{equation}
回放末值 $E_{144}$ 作为次日实际 $E_0$。

\subsubsection{参数取值与依据}

\begin{enumerate}
  \item \textbf{XGB 超参数。}树深 $6$、学习率 $0.05$、$200$ 棵树、行/列抽样 $0.8$、叶最小权重 $4$、$L_2$ 罚 $1$、种子 $42$；不按 $2$--$12$ 月电费交叉验证，避免用测试期信息选参。
  \item \textbf{光伏窗长 $W=7$。}见日前点预测：$1$ 月 $15$--$31$ 日同刻 MAE 在 $W=5\sim7$ 见底。
  \item \textbf{净分位 $\alpha=0.8$。}由报童临界覆盖率 $4\pi/(4\pi+\pi)=0.8$ 直接给定，残差取决策日前全部历史。
  \item \textbf{日末 $S^{*}=2400$ 与 $\rho=0.625$。}日末留电若回到中位 $6000\,\mathrm{kWh}$，当天谷电大量留给次日；若贴近下限 $1200\,\mathrm{kWh}$，次日早高峰只能靠购电，故 $S^{*}$ 存在折中。$\rho$ 过小会提前把次日要用的电放完，过大则日内缺口几乎全买紧急电。两项参数用 $1$ 月 $15$--$31$ 日做因果标定：以 $1200\,\mathrm{kWh}$ 为步长，在 $S^{*}\in\{2400,3600,4800,6000\}$、$\rho\in\{0.5,0.625\}$、$\alpha\in\{0.8,0.825\}$ 的 $16$ 个组合上，从储电量 $6000\,\mathrm{kWh}$ 起按当天真值反事实回放，并把 $1$ 月 $31$ 日 $24{:}00$ 的储电量折成现金计入总代价，最低的一组是 $S^{*}=2400$、$\rho=0.625$（图~\ref{fig:q2_jan_grid}）。
\end{enumerate}

\begin{figure}[!htbp]
  \centering
  \includegraphics[width=0.9\textwidth]{q2_january_sensitivity_grid.pdf}
  \caption{$1$ 月参数标定网格图}
  \label{fig:q2_jan_grid}
\end{figure}

由图~\ref{fig:q2_jan_grid} 可见：$\alpha=0.825$ 整面比 $\alpha=0.8$ 贵约 $2.7$ 万元，故不取更厚的分位。

每天 $0{:}00$ 重训 XGB，并求解一次含 $144\times5$ 个连续变量的日前线性规划，
$334$ 天均求得最优解且无约束校验错误。全年（$2025$ 年 $2$ 月 $1$ 日至 $12$ 月 $31$ 日）
总购电费为 $\textbf{13\,695\,403.61}$ 元，其中计划购电费 $13\,184\,604.50$ 元、
紧急购电费 $510\,799.11$ 元；紧急购电量 $87\,331.34$\,kWh，
分布在 $137$ 天共 $825$ 个时段；全年计划购电量 $21\,521\,947.29$\,kWh，
弃电量 $2\,472\,688.96$\,kWh。实际储电量始终位于
$[1200,10800]$\,kWh，最低为 $1200$\,kWh、最高为 $10800$\,kWh，
且无同时充放电。$12$ 月 $31$ 日 $24{:}00$ 储电量为 $2678.35$\,kWh。

全年日购电费呈明显的“冬高夏低”特征：$12$ 月日均 $5.46$ 万元最高，
$4$--$5$ 月为 $3.07$--$3.27$ 万元最低；$6$ 月因紧急购电出现全年最高单日费用，
当月紧急购电成本占全年紧急购电成本的 $37.2\%$。

题目指定日期的计划购电与储能调度见表~\ref{tab:q2_purchase}、
表~\ref{tab:q2_storage}。

% =========================================================
% 表6、表7并排
% 通过不同 arraystretch 使上下边界尽量同时对齐
% =========================================================

\begin{samepage}

\begin{table}[H]
  \centering

  % ==================== 左：表6 ====================
  \begin{minipage}[t]{0.49\textwidth}
    \vspace{0pt}
    \centering

    \caption{问题二：指定日期计划购电量及全天购电费}
    \label{tab:q2_purchase}

    \vspace{0.25em}

    \scriptsize
    \setlength{\tabcolsep}{1.2pt}

    \renewcommand{\arraystretch}{1.00}

    \begin{tabular*}{\linewidth}{
      @{\extracolsep{\fill}}
      ccccc
      @{}
    }
      \toprule
      时间段
        & 3.20
        & 6.21
        & 9.23
        & 12.21 \\
      \midrule

      10:00--10:10
        & 0.00
        & 0.00
        & 0.00
        & 0.00 \\

      12:00--12:10
        & 641.27
        & 0.00
        & 554.73
        & 1052.69 \\

      14:00--14:10
        & 0.00
        & 0.00
        & 0.00
        & 0.00 \\

      16:00--16:10
        & 466.08
        & 127.67
        & 475.19
        & 841.17 \\

      18:00--18:10
        & 743.26
        & 437.62
        & 806.06
        & 739.53 \\

      20:00--20:10
        & 0.00
        & 0.00
        & 0.00
        & 0.00 \\

      \midrule

      全天计划购电量（kWh）
        & 70698.38
        & 34922.73
        & 69093.42
        & 99124.96 \\

      全天计划购电费（元）
        & 43184.03
        & 20252.71
        & 42877.32
        & 64123.11 \\

      紧急购电费（元）
        & 0.00
        & 0.00
        & 1933.20
        & 0.00 \\

      全天总购电费（元）
        & 43184.03
        & 20252.71
        & 44810.52
        & 64123.11 \\

      \bottomrule
    \end{tabular*}

  \end{minipage}
  \hfill

  % ==================== 右：表7 ====================
  \begin{minipage}[t]{0.49\textwidth}
    \vspace{0pt}
    \centering

    \caption{问题二：指定日期储能充放电量及起止储电量}
    \label{tab:q2_storage}

    \vspace{0.25em}

    \scriptsize
    \setlength{\tabcolsep}{0.8pt}

    \renewcommand{\arraystretch}{1.10}

    \begin{tabular*}{\linewidth}{
      @{\extracolsep{\fill}}
      ccccc
      @{}
    }
      \toprule

      时间段
        & 3.20
        & 6.21
        & 9.23
        & 12.21 \\

      & \multicolumn{4}{c}{充电量 / 放电量（kWh）} \\

      \midrule

      0:00--4:00
        & 7718.29/49.42
        & 456.07/118.42
        & 7992.53/0.00
        & 7847.31/144.43 \\

      4:00--8:00
        & 969.00/4426.56
        & 542.27/1735.80
        & 722.97/5767.24
        & 1521.31/667.36 \\

      8:00--12:00
        & 5458.12/3265.56
        & 9472.90/0.00
        & 4774.60/2387.94
        & 5135.59/7756.25 \\

      12:00--16:00
        & 4053.33/266.64
        & 0.00/0.00
        & 5581.23/530.20
        & 7180.91/2220.11 \\

      16:00--20:00
        & 111.19/4859.30
        & 89.44/5164.11
        & 123.97/4997.03
        & 0.00/4779.93 \\

      20:00--24:00
        & 1815.79/3681.54
        & 2268.71/3017.50
        & 1948.30/3517.84
        & 1929.15/3622.79 \\

      \midrule

      $0{:}00$ 储电量
        & 3308.16
        & 3436.12
        & 2993.40
        & 3270.22 \\

      $24{:}00$ 储电量
        & 3033.50
        & 3831.65
        & 2911.26
        & 3199.88 \\

      \bottomrule
    \end{tabular*}

  \end{minipage}

\end{table}

\end{samepage}

\vspace{-0.4em}

\noindent

由表~\ref{tab:q2_purchase}、表~\ref{tab:q2_storage} 可见：
计划购电主要集中在光伏回落的 $16{:}00$、$18{:}00$ 时段，
$10{:}00$ 与 $20{:}00$ 则主要由富余光伏和储能放电覆盖；
储能总体呈“凌晨低价充电、早晚高峰放电”的运行规律。
四个指定日日末储电量均高于计划硬约束 $2400$\,kWh。
其中仅 $9$ 月 $23$ 日出现紧急购电，共 $288.88$\,kWh，
主要集中在晚间负荷高峰。

\end{samepage}


% =========================================================
% 指定日调度图
% =========================================================

指定日 2025-03-20 的 $144$ 个 $10\,\mathrm{min}$ 时段调度
如图~\ref{fig:q2-circular}。

\begin{samepage}

\begin{figure}[H]
  \centering
  \includegraphics[width=0.72\textwidth]{q2_circular_dispatch.pdf}
  \caption{指定日 (2025-03-20) 分时电价、净负载、充放电与储电量环形对照}
  \label{fig:q2-circular}
\end{figure}

\vspace{-0.5em}

\noindent

由图~\ref{fig:q2-circular} 可见，储能状态随全天电价和净负载滚动调整，
而非强制恢复至 $6000$\,kWh。该日凌晨低价时段集中充电，
早晚高峰集中放电：$0{:}00$--$4{:}00$ 充电 $7718.29$\,kWh，
$4{:}00$--$8{:}00$ 与 $16{:}00$--$20{:}00$ 分别放电
$4426.56$\,kWh 和 $4859.30$\,kWh。
SOC 始终处于 $[1200,10800]$\,kWh 范围内，
日初为 $3308.16$\,kWh、日末为 $3033.50$\,kWh。
受储能容量上限影响，当日产生弃电 $4.74$\,MWh，
但无紧急购电和缺供。该日全天计划购电 $70\,698.38$\,kWh，
总购电费为 $43\,184.03$ 元。

\end{samepage}


% =========================================================
% 调度环节对照
% =========================================================

\subsection{调度策略对照与效果分析}

在同一套点预测（负载扩展窗 XGBoost、光伏 $7$ 日同刻）
与 $2$ 月 $1$ 日储电量 $6000$\,kWh 下，
对不同分位口径及运行约束进行对照，结果见表~\ref{tab:q2_control}。

\begin{samepage}

\begin{table}[H]
  \centering
  \scriptsize
  \setlength{\tabcolsep}{3.2pt}
  \renewcommand{\arraystretch}{0.95}

  \caption{同一套点预测下不同调度策略的全年费用对照}
  \label{tab:q2_control}

  \begin{tabular}{lrrrr}
    \toprule
    策略
      & 计划购电费（元）
      & 紧急购电费（元）
      & 紧急电量（kWh）
      & 总购电费（元） \\
    \midrule

    无分位
      & 12124833.88
      & 3198303.43
      & 525856.18
      & 15323137.31 \\

    固定分量 $0.8/0.2$
      & 13369773.46
      & 395394.12
      & 66910.43
      & 13765167.58 \\

    净分位 $\alpha=0.8$
      & 13192124.98
      & 522427.80
      & 87329.02
      & 13714552.77 \\

    \midrule

    净分位 $+$ 日末 $2400$
      & 13185526.68
      & 513021.95
      & 82852.85
      & 13698548.63 \\

    净分位 $+$ 日末 $2400$ $+$ $\rho=0.625$
      & 13184604.50
      & 510799.11
      & 87331.34
      & \textbf{13695403.61} \\

    \bottomrule
  \end{tabular}
\end{table}

\vspace{-0.4em}

\noindent

由表~\ref{tab:q2_control} 可见，
在三种分位口径中，净分位策略较固定分量方案全年总费用降低约
$5.06$ 万元，较无分位方案降低约 $160.9$ 万元；
在净分位基础上加入日末 $2400$\,kWh 约束后，总费用进一步降低约
$1.6$ 万元，再加入 $\rho=0.625$ 后继续降低 $3145$ 元。
因此，最终采用“净分位 $+$ 日末 $2400$\,kWh $+$ $\rho=0.625$”
作为正式调度方案。回测区间逐日费用及分月分布见附录图~\ref{fig:app_q2_daily}、图~\ref{fig:app_q2_month}。

\end{samepage}
